% ============================================================
%  CAPÍTULO I: PLANTEAMIENTO DEL PROBLEMA Y JUSTIFICACIÓN
% ============================================================

\chapter{PLANTEAMIENTO DEL PROBLEMA Y JUSTIFICACIÓN}
\label{cap:problema}

% ─────────────────────────────────────────────────────────────
\section{Planteamiento del Problema}
\label{sec:planteamiento}

La seguridad del paciente constituye uno de los principales retos de los sistemas de salud contemporáneos. Según la Organización Mundial de la Salud \mbox{(OMS, 2019)}, los eventos adversos derivados de la atención sanitaria figuran entre las diez principales causas de muerte y discapacidad en el mundo, generando además costos sustanciales para los sistemas asistenciales. En el Perú, el Registro Nacional de Instituciones Prestadoras de Servicios de Salud (RENIPRESS) registra 26{,}663 establecimientos activos que atienden a más de 33 millones de personas, distribuidos en doce tipos institucionales: sector privado (16{,}685 IPRESS), Gobierno Regional (8{,}380), MINSA (510), EsSalud (408), Sanidades de las Fuerzas Armadas y Policiales, SISOL, gobiernos locales y otros. EsSalud, que concentra la mayor densidad de hospitalizaciones de tercer nivel, aprobó en 2021 la Directiva GG-ESSALUD «Registro, Notificación y Gestión de los Eventos Relacionados con la Seguridad del Paciente (ERSP)», la cual incorpora explícitamente la Matriz de Priorización del Modelo de Gestión Moderna de los Servicios de Salud (GEMSES) como herramienta normativa para clasificar el nivel de riesgo y de gestión de los eventos.

Sin embargo, el proceso de identificación y priorización depende todavía del registro voluntario de los profesionales de la salud y del análisis manual de las descripciones libres consignadas en las notas clínicas. Esta dependencia genera tres obstáculos persistentes: (i)~la subnotificación voluntaria, asociada al temor de repercusiones laborales; (ii)~la dependencia de descripciones en lenguaje natural que dificultan su codificación, agregación y análisis sistemático; y (iii)~la falta de mecanismos automatizados que prioricen los eventos de mayor riesgo para escalarlos al nivel de gestión adecuado en tiempo oportuno.

En el caso específico de EsSalud, la Directiva GG-ESSALUD-2021 establece un flujo formal de Registro $\rightarrow$ Notificación $\rightarrow$ Validación $\rightarrow$ Codificación $\rightarrow$ Priorización $\rightarrow$ Mitigación $\rightarrow$ Análisis de Causa Raíz $\rightarrow$ Plan de Acción. Aunque la directiva normaliza la Matriz de Priorización GEMSES como instrumento oficial, su aplicación efectiva requiere que personal de la Oficina de Gestión de la Calidad y Humanización procese manualmente cada notificación, lo que constituye un cuello de botella estructural que hace el proceso lento, costoso y vulnerable a la heterogeneidad en la interpretación.

La presente investigación propone aprovechar los avances recientes en Procesamiento de Lenguaje Natural (NLP) y aprendizaje automático para automatizar la detección de 188 de los 231 eventos adversos del Anexo 02 de la Directiva GG-ESSALUD-2021 en notas clínicas no estructuradas, y la posterior aplicación de la Matriz de Priorización GEMSES. Dada la ausencia de acceso a registros institucionales peruanos, el estudio se desarrollará sobre el \textit{dataset} abierto MIMIC-IV (\textit{Medical Information Mart for Intensive Care}, versión~IV). El estudio se concibe como un \textbf{demostrador técnico de viabilidad}: demuestra que la automatización GEMSES es realizable con datos clínicos reales, generando un protocolo documentado y replicable que las instituciones peruanas ---cuyo sistema nacional aún carece de soporte normativo e infraestructura para esta automatización--- podrán adoptar cuando se habiliten los accesos correspondientes.

% ─────────────────────────────────────────────────────────────
\section{Formulación del Problema}
\label{sec:formulacion}

La gestión de eventos adversos en los servicios de salud presenta brechas críticas en la detección oportuna, la codificación consistente y la priorización basada en evidencia. Los modelos manuales dependen del criterio subjetivo del profesional y generan subnotificación sistemática. Más aún, incluso cuando un evento se reporta, la cadena de gestión presenta cuatro puntos de quiebre sucesivos: (i)~el evento se registra pero \textit{no se analiza}, por saturación del personal de calidad; (ii)~se analiza pero \textit{no se prioriza}, por ausencia de herramientas de clasificación estandarizada; (iii)~se prioriza pero \textit{no se asigna un responsable} con plazo definido; y (iv)~se gestiona puntualmente pero \textit{no existe un indicador permanente de control de calidad} que mida la recurrencia del evento. Este patrón impide que los sistemas de salud aprendan de sus errores y escalen la vigilancia frente al volumen de notas clínicas generadas en hospitales de tercer nivel.

En consecuencia, el problema central se formula mediante la siguiente pregunta de investigación:

\begin{quote}
\textit{¿En qué medida un pipeline integrado de procesamiento de lenguaje natural y aprendizaje automático supervisado, aplicado sobre notas clínicas no estructuradas del \textit{dataset} MIMIC-IV, permite detectar eventos adversos y aplicar la Matriz de Priorización del Modelo GEMSES ---incluyendo la asignación automática del Nivel de Gestión y del responsable institucional--- con una concordancia significativamente mayor a la clasificación manual experta tradicional?}
\end{quote}

% ─────────────────────────────────────────────────────────────
\section{Justificación}
\label{sec:justificacion}

La presente investigación se justifica desde cuatro dimensiones complementarias:

% - - - - - - - - - - - - - - - - - - - - - - - - - - - - - -
\subsection{Justificación Teórica}
\label{subsec:justif_teorica}

El campo de la detección automática de eventos adversos en texto clínico ha avanzado sustancialmente en los últimos años, con trabajos basados en sistemas de supervisión débil (Ratner et al., 2017), modelos de transformers biomédicos (Lee et al., 2020; Alsentzer et al., 2019) y técnicas de negación clínica como NegEx (Chapman et al., 2001). Sin embargo, ningún trabajo previo ha articulado estos enfoques con la Matriz de Priorización del Modelo GEMSES ni con la taxonomía de 231 eventos del Anexo 02 de la Directiva GG-ESSALUD-2021. Esta investigación cierra esa brecha teórica, aportando un marco de operacionalización que vincula las dimensiones de Tiempo, Calidad, Costos y Satisfacción del Modelo GEMSES con variables derivables de datos clínicos estructurados y no estructurados. Los pesos de impacto por dimensión ya están precalculados en el Anexo~02 para cada uno de los 231 eventos; esta investigación los operacionaliza computacionalmente.

Adicionalmente, el estudio contribuye al debate metodológico sobre la transferibilidad de modelos de NLP entrenados en inglés clínico hacia sistemas sanitarios de habla hispana, problema de creciente relevancia en la literatura de informática biomédica internacional.

% - - - - - - - - - - - - - - - - - - - - - - - - - - - - - -
\subsection{Justificación Práctica}
\label{subsec:justif_practica}

Desde una perspectiva práctica, el sistema propuesto responde a una necesidad institucional concreta: el cuello de botella operativo que genera el procesamiento manual de notificaciones de eventos adversos en EsSalud. El \textit{pipeline} automatizado permitiría procesar el volumen completo de notas clínicas de alta (331{,}793 notas en MIMIC-IV; equivalente a decenas de miles en EsSalud) sin incremento proporcional de recursos humanos especializados, reduciendo el tiempo entre la ocurrencia del evento y su priorización para la gestión correctiva.

El modelo validado sobre MIMIC-IV generará un protocolo de implementación documentado y replicable, que podrá ser adaptado por la Oficina de Gestión de la Calidad y Humanización de EsSalud ---o cualquier institución que opere bajo estándares GEMSES--- cuando se habilite el acceso a datos clínicos institucionales para investigación. El impacto práctico esperado incluye: reducción de la subnotificación, mayor consistencia en la codificación, y escalamiento de la capacidad de vigilancia de la seguridad del paciente.

% - - - - - - - - - - - - - - - - - - - - - - - - - - - - - -
\subsection{Justificación Metodológica}
\label{subsec:justif_metodologica}

El uso de MIMIC-IV como corpus de validación garantiza la reproducibilidad del estudio: es un \textit{dataset} de acceso controlado pero abierto (PhysioNet \textit{credentialing}), con 145{,}914 pacientes únicos, en formato estándar internacional, y ampliamente utilizado como \textit{benchmark} en investigación de NLP clínico. La elección de este \textit{dataset} ---en lugar de datos institucionales peruanos no accesibles--- no es una limitación sino una decisión metodológica deliberada que garantiza la trazabilidad, la posibilidad de revisión por pares y la comparabilidad con el estado del arte internacional.

El diseño del \textit{pipeline} combina supervisión débil para etiquetado a escala, detección de negación con NegEx, mapeo a códigos ICD-10, y comparación de múltiples clasificadores (TF-IDF + regresión logística, LinearSVC, modelos basados en \textit{transformers}). La validación mediante panel experto ($n = 70$ notas) con métricas de concordancia (kappa de Cohen) proporciona evidencia estadística rigurosa sobre la capacidad del sistema, con un IC$_{95\%}$ Wilson calculado correctamente para muestras pequeñas.

% - - - - - - - - - - - - - - - - - - - - - - - - - - - - - -
\subsection{Justificación Social}
\label{subsec:justif_social}

La OMS estima que en países de ingresos medios y bajos, la carga de eventos adversos evitables equivale a 134 millones de episodios anuales con 2.6 millones de muertes asociadas (OMS, 2019). Perú, como país de ingreso medio-alto con un sistema de seguridad social en expansión, enfrenta este problema de manera aguda: la red asistencial de EsSalud atiende a más de 13 millones de asegurados con capacidad limitada de vigilancia activa de seguridad, que es detectada cuando el paciente denuncia; SUSALUD ha multado a instituciones por eventos de seguridad del paciente evitables por un valor de S/~8{,}482{,}906.50 en el año 2025.

Un sistema automatizado de detección y priorización contribuye directamente a la protección de pacientes vulnerables, a la rendición de cuentas institucional y al aprendizaje organizacional. La generación de alertas tempranas sobre eventos de nivel Rojo (alta prioridad GEMSES) permitiría activar protocolos de mitigación antes de que el daño escale, con un impacto social concreto y medible en términos de morbimortalidad prevenible y sostenibilidad financiera.
